\documentclass[../DoAn.tex]{subfiles}
\begin{document}

\begin{center}
    \Large{\textbf{TÓM TẮT NỘI DUNG ĐỒ ÁN}}\\
\end{center}
\vspace{1cm}
Đồ án xây dựng ThePawsome, một hệ thống thương mại điện tử chuyên biệt cho sản phẩm thú cưng, kết hợp các nghiệp vụ mua bán trực tuyến với trợ lý AI tư vấn bằng tiếng Việt. Bài toán xuất phát từ thực tế người nuôi thú cưng thường phải tự lựa chọn thức ăn, phụ kiện và kiến thức chăm sóc từ nhiều nguồn phân tán, trong khi mỗi thú cưng có loài, độ tuổi, cân nặng, tình trạng sức khỏe và dị ứng khác nhau. Các website bán hàng thông thường đáp ứng tốt nhu cầu trưng bày sản phẩm và đặt hàng, nhưng còn hạn chế ở khả năng cá nhân hóa tư vấn, kiểm soát chất lượng tri thức cộng đồng và bảo đảm giao dịch an toàn khi nhiều người cùng mua một biến thể sản phẩm.

Giải pháp của đồ án là thiết kế hệ thống web theo kiến trúc tách frontend/backend. Backend sử dụng FastAPI, SQLAlchemy bất đồng bộ, PostgreSQL tích hợp pgvector và Redis; frontend sử dụng Next.js App Router, TypeScript, Tailwind CSS, TanStack Query và Zustand. Trợ lý Catbot được triển khai bằng LangGraph tool-calling có lớp planner LLM tạo kế hoạch tool-call trước khi agent trả lời, truy xuất sản phẩm, kho kiến thức, thảo luận forum và hồ sơ thú cưng từ CSDL để sinh câu trả lời có căn cứ, đồng thời có cơ chế an toàn AI và chuyển giao cho nhân viên hỗ trợ. Hệ thống thương mại điện tử hỗ trợ catalog có biến thể, giỏ hàng, guest checkout, mã giảm giá, đánh giá sau mua, đổi trả, thanh toán QR qua SePay, giữ kho tạm thời và xử lý idempotency để hạn chế tạo đơn trùng. Kết quả đạt được là một monorepo có đầy đủ các phân hệ khách hàng, quản trị, diễn đàn, AI/RAG và kiểm thử backend trọng tâm, làm nền tảng để tiếp tục triển khai thực nghiệm với dữ liệu thật.
\begin{flushright}
Sinh viên thực hiện\\
\begin{tabular}{@{}c@{}}
\textit{(Ký và ghi rõ họ tên)}
\end{tabular}
\end{flushright}

\end{document}
